% Independent audit addendum; not a proof of RH.
% Assumes the continuation report's exact h_a and finite-window compactness.
\section{An exact all-window boundedness criterion}
Fix $a>0$, and let $h_a$ be the real, locally integrable arithmetic kernel
defined above, supported in $[1,\infty)$.  Define the finite-window operators
directly by
\[
 (H_{a,L}f)(x)=\int_0^L h_a(xy)f(y)\,dy,
 \qquad f\in L^2((0,L),dx),\quad L>1.
\]
Their boundedness and compactness do not assume the existence of a bounded
global operator.

\begin{theorem}[All-window alternative]
For each fixed $a>0$, the following are equivalent:
\begin{enumerate}
\item $\sup_{L>1}\|H_{a,L}\|<\infty$;
\item $\|H_{a,L}\|\leq1$ for every $L>1$;
\item $R_a(w)=\xi(1/2+w-a)/\xi(1/2+w+a)$ is holomorphic and bounded by one
      on $\operatorname{Re}w>0$.
\end{enumerate}
The function $L\mapsto\|H_{a,L}\|$ is nondecreasing.  Its limit as
$L\to\infty$ is $1$ when these conditions hold and $+\infty$ otherwise.
\end{theorem}

\begin{proof}
For $1<L_1<L_2$, the first operator is the orthogonal compression of the
second, so its norm is no larger.  Suppose condition (1) holds with bound
$C$.  For compactly supported $f\in L^2((0,\infty),dx)$ and $L$ beyond
its support, $H_{a,L}f$ is the restriction to $(0,L)$ of a single locally
defined function.  The uniform norm bound and monotone convergence give
a global $L^2$ function $H_af$ with $\|H_af\|\leq C\|f\|$.  Density
therefore yields a bounded global extension $H_a$.

The operator
\[
 (Jf)(x)=x^{-1}f(x^{-1})
\]
is unitary on $L^2((0,\infty),dx)$.  Substituting $y=t^{-1}$ yields
\[
 (H_aJf)(x)=\int_0^\infty h_a(x/t)f(t)\,\frac{dt}{t}.
\]
For $f\in L^2((1,\infty),dx)$, support of $h_a$ restricts the integral
to $1\leq t\leq x$.  Local $L^1$ convolution bounds in logarithmic
coordinates, followed by truncation of $f$, identify this integral with
the bounded extension $H_aJf$ for every such $f$.  Thus condition (3) of
Suzuki's Theorem 2.2 \cite{Suzuki2016} holds, with $\omega=a$.  It follows
that $\Theta_a(z)=R_a(-iz)$ is inner in the upper half-plane, giving (3).

Conversely, (3), the reflection symmetries of $\xi$, and Suzuki's
Lemma 4.1 identify the global arithmetic kernel operator $H_a$ as an
isometry.  Consequently all finite compressions have norm at most one,
giving (2).  Condition (2) plainly gives (1).

When these conditions hold, the compressions converge strongly to $H_a$
as $L\to\infty$, so their norm supremum is $\|H_a\|=1$.  Otherwise (1)
fails and monotonicity implies that their norms tend to infinity.
\end{proof}

\begin{corollary}[The required small-width estimate]
Let $a_n>0$ tend to zero.  If for each $n$ there exists a finite constant
$C_n$ such that
\[
 \|H_{a_n,L}\|\leq C_n\qquad\hbox{for every }L>1,
\]
then RH holds.  No uniformity of $C_n$ in $n$, and no strictly positive
gap below one uniform in $L$, is required.
\end{corollary}

\begin{proof}
The theorem makes every $R_{a_n}$ holomorphic in the right half-plane.
If $X(w)=\xi(1/2+w)$ had a zero $\lambda$ with positive real part,
then for sufficiently small $a_n$ the point $w=\lambda-a_n$ lies in that
half-plane.  The denominator of $R_{a_n}$ vanishes there, while its
numerator $X(\lambda-2a_n)$ does not, by isolation of the zero $\lambda$.
This is a pole, a contradiction.  Evenness of $X$ excludes zeros with
negative real part as well.
\end{proof}

\paragraph{Scope.}
The all-window bounds in this corollary are not established here.
The finite certificate $a=1/4$, $1<L\leq3/2$ does not provide them.
This criterion identifies the remaining arithmetic estimate; it neither
derives the scalar factor from twistor geometry nor constructs a
canonical-system Hamiltonian.  The latter requires separate regularity
and normalization arguments.

% Bibliography entry for integration:
% \bibitem{Suzuki2016} M. Suzuki, A canonical system of differential equations
% arising from the Riemann zeta-function, arXiv:1204.1827, revised version,
% Theorem 2.2 and Lemma 4.1. https://arxiv.org/abs/1204.1827
